\section{Introduction}
\label{sec:intro}

Electroencephalography (EEG) foundation models --- transformer encoders
pretrained on large unlabelled EEG corpora via masked reconstruction or
contrastive objectives --- have been proposed as a general-purpose feature
extractor for clinical and affective decoding
\citep{labram2024,cbramod2025,reve2025}. On resting-state stress
detection, a high-impact clinical target where self-report instruments
such as the Depression--Anxiety--Stress Scale (DASS) provide trustworthy
labels and acquisition is cheap, recent work reports balanced accuracy
(BA) above 90\%: Wang et al.\ \citep{wang2025stress} fine-tune LaBraM on
the Komarov et al.\ Stress dataset \citep{komarov2020stress} and achieve
$\mathrm{BA} = 0.9047$. Taken at face value, this suggests near
clinical-grade stress screening from a five-minute eyes-open recording.

We attempted to reproduce this result and found a 30+ percentage-point
gap that does not reflect model quality. The key divergence lies in
cross-validation protocol: the reported headline uses \emph{trial-level}
splits, in which different recordings from the same subject may appear
in both the training and test partitions. Because EEG exhibits strong
individual stability, trial-level splits implicitly reward subject
re-identification rather than stress-state classification. Under
\emph{subject-level} StratifiedGroupKFold with honest per-recording DASS
labels, cuDNN-deterministic training, and a multi-seed protocol, the
same three foundation models --- LaBraM, CBraMod, and REVE --- yield
frozen linear-probe BA between $0.45$ and $0.60$ and best-hyperparameter
fine-tuned BA between $0.52$ and $0.58$ on Stress.

The gap between trial-level and subject-level evaluation is not a novel
observation in itself. Roy et al.\ \citep{roy2019review} documented
subject leakage as a pervasive issue across the EEG deep learning
literature. Three independent 2025--2026 benchmarks ---
Brain4FMs~\citep{brain4fms2026}, EEG-FM-Bench~\citep{eegfmbench2025}, and
AdaBrain-Bench~\citep{adabrain2025} --- have since converged on the
finding that EEG foundation models fail to outperform classical
baselines on cross-subject affective and cognitive tasks. Our
contribution is not to re-validate this failure, but to provide the
\emph{mechanistic} explanation these benchmarks lack.

We organise the diagnosis around three questions.
\textbf{(Q1) What do frozen EEG foundation-model representations
encode?} We show via representational similarity analysis (RSA) and
mixed-effects variance decomposition that subject identity dominates the
learned feature space across twelve model$\times$dataset combinations,
yet a low signal-to-noise linearly separable label subspace survives ---
one that classical band-power features cannot access. \textbf{(Q2) When
does fine-tuning help, when does it hurt, and why?} Across three
foundation models and four datasets (Stress, ADFTD, TDBRAIN, EEGMAT) we
find that fine-tuning direction is a model$\times$dataset interaction:
LaBraM injects label signal on ADFTD but erodes it on TDBRAIN, while
REVE does the opposite on the same labels. Within-subject contrast
strength, not the presence of a within-subject design, determines
whether fine-tuning rescues a given dataset. \textbf{(Q3) Can
small-sample EEG benchmarks support claims of foundation-model
superiority?} Using the 70-recording Stress dataset as a cross-pillar
stress test, we present five independent lines of evidence
(cuDNN-induced single-seed variance, permutation-null indistinguishability,
a 2017-vintage ShallowConvNet matching foundation-model performance
from scratch, cross-model divergence in band-stop ablation, and
within-subject longitudinal failure) that converge on a single
conclusion: this regime lies below the statistical power floor required
to support any foundation-model scientific claim.

The rest of the paper is organised as follows. Section~\ref{sec:related}
situates our work among EEG foundation models, the Stress dataset, and
the 2025--2026 benchmark literature. Section~\ref{sec:methods} describes
our dataset preparation, cross-validation protocol, per-model input
normalisation (a silent failure mode if mishandled), multi-seed
determinism, and statistical reporting conventions. Section~\ref{sec:results}
presents the three-pillar diagnosis and the neuroscience interpretability
triad, culminating in the Stress power-floor case study
(Section~\ref{sec:power_floor}). Section~\ref{sec:discussion} offers
hypotheses for the model$\times$dataset interaction and field-level
recommendations; Section~\ref{sec:conclusion} summarises.
